\section{InnerProduct root scalar consumer frontier}
\label{sec:innerproduct-root-scalar-consumer-frontier}

\begin{theorem}[InnerProduct root scalar projection transport]
\label{thm:innerproduct-root-scalar-projection-transport}
In the root obligation basis of
\autoref{def:innerproduct-root-unblock-obligation-basis}, the scalar endpoint
of the form transports through the retained scalar classifier. Concretely, if
$x\sim_Vx'$ and $y\sim_Vy'$ are vector-classifier witnesses for carried
endpoints $x,y,x',y':C_V$, then the projected form endpoints satisfy
$$
\langle x,y\rangle_{\InnerProductUp}
\sim_{\mathsf{scal}}
\langle x',y'\rangle_{\InnerProductUp}.
$$
This projection uses only the root form row, the two vector-classifier
witnesses, and the retained scalar classifier.
\end{theorem}
\begin{proof}
Read the form endpoint and both classifiers from
\autoref{def:innerproduct-root-unblock-obligation-basis}. The basis exposes
the vector classifier $\sim_V$, the scalar classifier
$\sim_{\mathsf{scal}}$, and the form endpoint row, with no additional scalar
equality.

Apply \autoref{thm:innerproduct-classifier-transport-row} to the two
vector-classifier witnesses. Its conclusion is exactly a retained
scalar-classifier witness between the two form endpoints. Thus form endpoint
transport is a scalar projection of the root basis data.
\end{proof}

\begin{theorem}[InnerProduct root diagonal zero exactness]
\label{thm:innerproduct-root-diagonal-zero-exactness}
For every carried vector endpoint $x:C_V$ in the root obligation basis of
\autoref{def:innerproduct-root-unblock-obligation-basis}, the diagonal
self-value $\langle x,x\rangle_{\InnerProductUp}$ is classified as the scalar
zero endpoint if and only if $x\sim_V0_V$. The equivalence is a
$\mathsf{BHist}$ row statement: scalar-zero classification of the diagonal is
exactly the displayed vector-zero classifier class.
\end{theorem}
\begin{proof}
First read the diagonal endpoint from the form-law surface of
\autoref{thm:innerproduct-root-unblock-form-laws}. This places
$\langle x,x\rangle_{\InnerProductUp}$ inside the displayed scalar endpoint
source and keeps its classifier equal to $\sim_{\mathsf{scal}}$.

For the forward direction, assume that the diagonal endpoint is classified as
scalar zero. Apply
\autoref{thm:innerproduct-positive-definite-obligation}. Its exactness clause
converts this scalar-zero classification into the vector-classifier witness
$x\sim_V0_V$.

For the reverse direction, assume $x\sim_V0_V$. The converse exactness clause
of \autoref{thm:innerproduct-positive-definite-obligation}, together with the
root form-law unblock row, returns the scalar-zero classification of
$\langle x,x\rangle_{\InnerProductUp}$. Hence the diagonal zero test is exact
at the root surface.
\end{proof}

\begin{theorem}[InnerProduct root Norm-Hilbert consumer frontier]
\label{thm:innerproduct-root-norm-hilbert-consumer-frontier}
The $\NormUp$ and $\HilbertUp$ consumers of an $\InnerProductUp$ root
obligation basis may read only the following scalar-consumer frontier:
the norm-squared diagonal endpoint $\|x\|_{I}^{2}=\langle x,x\rangle$,
the diagonal zero exactness row, the orthogonality row, the carrier and
classifier rows, and the visible ledger rows. They receive no completeness,
topology, host Hilbert-space carrier, or Cauchy-limit selector from
$\InnerProductUp$.
\end{theorem}
\begin{proof}
For $\NormUp$, use \autoref{def:innerproduct-norm-metric-seed-source}. It
defines the norm-squared seed from the diagonal form endpoint and lists the
positive-definite exactness, classifier-transport, and ledger rows needed to
consume that endpoint.

Apply \autoref{thm:innerproduct-root-diagonal-zero-exactness} to identify
the zero test of the norm-squared diagonal with the displayed vector-zero
classifier class. Orthogonality is read from
\autoref{def:innerproduct-orthogonality-row}, so it is also a scalar-zero
classification of a displayed form endpoint.

For $\HilbertUp$, apply
\autoref{thm:innerproduct-root-threshold-orthogonality-consumption}. It
separates the inner-product carrier, classifier, form, transport,
orthogonality, and ledger rows from Banach metric and completeness data.

Finally use \autoref{thm:innerproduct-root-unblock-downstream-threshold}.
The stated root threshold contains no completeness field, topology field,
host Hilbert-space carrier, or Cauchy-limit selector. Therefore the
$\NormUp$ and $\HilbertUp$ consumer frontier is exactly the listed
scalar-consumer surface.
\end{proof}
